\section{The local exchange-rate bound and the preorder on the Pauli-plus-erasure class}\label{sec:exchange}

The main line of this section: \emph{discarding an erasure flag is a free operation} (Corollary~\ref{cor:discard}), which gives a rigorous upper bound on the exchange rate (Theorem~\ref{thm:rate}),
and in the simplest subclass $\Nn_{\rm cc}(p,e)$ we completely characterize the preorder induced by local free operations (Theorem~\ref{thm:preorder}).
The gap between the true exchange rate and the bound is one of the core research questions of the project (\S\ref{sec:gap}).

\subsection{Partial flag discarding and the exchange-rate bound}

\begin{theorem}[Partial flag-discarding inequality]\label{thm:partial}
Let $0\le p\le\tfrac34$, $0\le e_0<1$, $0\le\delta\le1-e_0$. For every code distance $d$,
\[
\err_d\big(\Nn_{\rm cc}(p,\,e_0+\delta)\big)\;\le\;\err_d\big(\Nn_{\rm cc}(p'',\,e_0)\big),\qquad
p''=p+\delta\,\frac{\tfrac34-p}{1-e_0}.
\]
Consequently $\alpha_\pm(p,e_0+\delta)\ge\alpha_\pm(p'',e_0)$.\hfill\st{PROVED}
\end{theorem}
\begin{proof}
Apply Lemma~\ref{lem:basic-ops}(c) with $e=e_0+\delta$ and $\kappa=\delta/(e_0+\delta)$ (trivial if $e_0+\delta=0$).
Then $e'=e(1-\kappa)=e_0$ and $e\kappa=\delta$, so
\[
p'=\frac{(1-e_0-\delta)p+\tfrac34\delta}{1-e_0}=p+\delta\frac{\tfrac34-p}{1-e_0}=p''.
\]
This operation is (F1$'$), so Theorem~\ref{thm:monotone} gives $\err(p,e_0+\delta)\le\err(p'',e_0)$;
the exponent version follows from Corollary~\ref{cor:alpha}.
\end{proof}

\begin{theorem}[Exchange-rate bound]\label{thm:rate}
Let $c(p_0,e_0)=\dfrac{3/4-p_0}{1-e_0}$.

(a) If the right derivatives of $\err_d$ with respect to $p$ and $e$ exist at $(p_0,e_0)$, then
\[\partial_e\err_d\le c\,\partial_p\err_d,\]
and when $\partial_p\err_d>0$, $R^{(d)}:=\partial_e\err_d/\partial_p\err_d\le c(p_0,e_0)$.

(b) Suppose the limit $\alpha=\alpha_-=\alpha_+$ exists, has right derivatives with respect to $p$ and $e$ at $(p_0,e_0)$, and $\partial_p\alpha<0$. Then the local exchange rate satisfies
\[
R_{e\to p}(p_0,e_0)=-\frac{\partial_e\alpha}{\partial_p\alpha}\;\le\;\frac{3/4-p_0}{1-e_0}.
\]
\hfill\st{PROVED}
\end{theorem}
\begin{proof}
In Theorem~\ref{thm:partial}, $p''=p_0+c\,\delta$ is linear in $\delta$. (a): subtract $\err_d(p_0,e_0)$ from both sides and divide by $\delta\downarrow0$;
the left side tends to $\partial_e\err_d$ and the right side to $c\,\partial_p\err_d$. (b): the same for $\alpha$ gives
$\partial_e\alpha\ge c\,\partial_p\alpha$; since $-\partial_p\alpha>0$, multiplying by $-1$ gives
$-\partial_e\alpha\le c(-\partial_p\alpha)$, and dividing by $-\partial_p\alpha$ gives the claim.
\end{proof}

\paragraph{Constitutional reading.}
$R_{e\to p}$ here is the \emph{marginal rate of a monotone} (\S\ref{sec:rates}); by Proposition~\ref{prop:marginal} every monotone of the local free order has marginal rate in $[0,c]$, and $c$ is attained by the free-order monotone $1-\mu$, so Theorem~\ref{thm:rate} cannot be improved using the free operations alone.

\begin{remark}[Relation to the heuristic ``$R\le3/4$'']\label{rem:34}
``Adding $\delta e$ of erasure $\to$ adding $\tfrac34\delta e$ of depolarizing noise'' is the value of $c$ at $(p_0,e_0)=(0,0)$. The exact $c$ has two corrections:
(i) the numerator $\tfrac34-p_0$: a newly erased qubit already carried depolarizing noise of probability $p_0$, so the net increase of the nontrivial-Pauli probability is $\tfrac34-p_0$;
(ii) the denominator $1-e_0$: $p$ is the conditional rate on \emph{non-erased qubits}, and the overall nontrivial-Pauli probability is $(1-e)p+\tfrac34 e$.

Hence $c\le\tfrac34$ if and only if $p_0\ge\tfrac34e_0$. At $e_0=0$ we always have $c=\tfrac34-p_0<\tfrac34$;
but at the work point $(p_0,e_0)=(0.02,0.05)$ suggested in the README, $c=0.768>\tfrac34$.
So ``$R\le3/4$'' is a theorem only in the region $p_0\ge\tfrac34e_0$; in general one should use $c(p_0,e_0)$.
\end{remark}

\begin{corollary}[Integral form: the ``Pauli equivalent'' of erasure]\label{cor:worth}
For every $(p,e)$ define the equivalent $p_{\rm eq}^{(d)}(p,e):=\inf\{p'\in[0,\tfrac34]:\err_d(p',0)\ge\err_d(p,e)\}$. Then
\[
p_{\rm eq}^{(d)}(p,e)\;\le\;p+e\,(\tfrac34-p),
\]
which is Theorem~\ref{thm:partial} with $e_0=0$ and $\delta=e$. If $\err_d(\cdot,0)$ is strictly increasing there is also the lower bound $p_{\rm eq}^{(d)}\ge p$
(monotonicity under superposed erasure, Proposition~\ref{prop:rate-mono}); the Pauli equivalent of the erasure is then sandwiched between $0$ and $e(\tfrac34-p)$,
that is, $0\le R_{\rm eff}\le\tfrac34-p$.\hfill\st{PROVED}
\end{corollary}

\begin{observation}[{Sharpness of the bound: the $[[5,1,3]]$ code attains equality at $e_0=0$}]\label{obs:sharp}
\texttt{theory/checks/exact\_small\_codes.py} enumerates all $4^n$ Pauli errors exactly for three codes with $n\le9$, computes the maximum-likelihood
$\err(p,e)$ in each case, and verifies Theorem~\ref{thm:partial} (all checks pass; see \S\ref{sec:map}). In particular:
\begin{itemize}
\item For the five-qubit code $[[5,1,3]]$ at $e_0=0$ and $p\in\{0.005,0.02,0.06,0.1,0.2,0.3,0.45\}$,
$R^{(d)}=\partial_e\err/\partial_p\err$ and $c=\tfrac34-p$ are \emph{equal} to 8 significant digits.
\item For the Steane code $R^{(d)}/c$ lies between $0.67$ and $0.86$; for the $d=3$ rotated surface code at $p=0.02$, $e_0=0$ one finds $R^{(3)}\approx0.54$ while $c=0.73$.
\end{itemize}
Explanation (unproved): for a perfect code, flagging a single qubit does not change the ML decision on any syndrome, so $\err$ is linear along the segment ``replace the law of qubit $j$ from
$\Dep_p$ by $\Dep_{3/4}$'', and the first-order coefficient is exactly $(\tfrac34-p)\partial_{p_j}\err$.
This shows that the inequality $\partial_e\err_d\le c\,\partial_p\err_d$ cannot be improved uniformly over all codes; the strict gap on the surface code is \emph{code dependent}.
\hfill\st{NUMERICAL}
\end{observation}

\subsection{Complete characterization of the local-free-operation preorder on the code-capacity class}

\begin{definition}[Local free operations]\label{def:local}
On the i.i.d.\ class, a \emph{local free operation} acts qubit by qubit independently: for each qubit the processor draws $(n,f')$ from a Markov kernel
$Q(n,f'\mid f)$ (with $f\in\{0,1\}$ the input flag, $n\in\F^2$ the superposed Pauli, and $f'\in\{0,1\}$ the output flag),
and sets $E'_j=E_j+n$, $F'_j=f'$. ($n$ may depend on $f$, $f'$, and the processor's randomness, but not on $E_j$; this is the general local form of (F1$'$), (F2), and (F2$'$).)
\end{definition}

\begin{theorem}[Preorder theorem: the Pauli-plus-erasure class]\label{thm:preorder}
Let $p,p'\in[0,\tfrac34)$ and $e,e'\in[0,1)$. The following are equivalent:
\begin{enumerate}[(i)]
\item there is a local free operation turning $\Nn_{\rm cc}(p,e)$ into $\Nn_{\rm cc}(p',e')$;
\item either $e'\ge e$ and $\lambda(p')\le\lambda(p)$ (that is, $p'\ge p$); or $e'<e$ and $\mu(p',e')\le\mu(p,e)$,
that is, $(1-e')\lambda(p')\le(1-e)\lambda(p)$.
\end{enumerate}
In particular, (i) implies $\err_d(p',e')\ge\err_d(p,e)$ for every code and every $d$.\hfill\st{PROVED}
\end{theorem}
\begin{proof}
\emph{(ii)$\Rightarrow$(i).} Case $e'\ge e$, $p'\ge p$: first superpose erasure by Lemma~\ref{lem:basic-ops}(b) to take $e\to e'$,
then superpose depolarizing noise by (a) to take $\lambda\to\lambda(p')$ (this needs $\lambda(q)=\lambda(p')/\lambda(p)\in[0,1]$).
Case $e'<e$, $\mu'\le\mu$: by (c) take $\kappa=1-e'/e$, which gives erasure rate $e'$ and
$\lambda_c=(1-e)\lambda(p)/(1-e')$. The condition $\mu'\le\mu$ is exactly $\lambda(p')\le\lambda_c$, and then (a) superposes the remaining depolarizing noise.

\emph{(i)$\Rightarrow$(ii).} For $P\in\F^2\setminus\{0\}$ write $\hat\nu(P)=\sum_a\nu(a)(-1)^{\langle P,a\rangle}$.
Input: with flag 1 (probability $e$) $E$ is uniform and $\hat E(P)=0$; with flag 0 (probability $1-e$) $\hat E(P)=\lambda:=\lambda(p)$.
Let $q_{f,f'}(n)=Q(n,f'\mid f)$ (a sub-probability measure) with total mass $q_{f,f'}$. The Fourier coefficient of the sub-measure with output flag $1$ is
$(1-e)\lambda\,\hat q_{0,1}(P)$ (the part coming from input flag 1 is uniform and has no nonzero modes).
The output class requires ``$E'$ uniform given flag 1'', so $\lambda\hat q_{0,1}(P)=0$ for all $P\ne0$; since $\lambda>0$, $q_{0,1}$ is a uniform measure.
Let $s_0=\sum_n q_{0,0}(n)$; then $\sum q_{0,1}=1-s_0$, and the probability of output flag $0$ is
\[
m':=1-e'=e\,y+(1-e)s_0,\qquad y:=\textstyle\sum_nq_{1,0}(n)\in[0,1].
\]
The Fourier coefficient at $P\neq0$ of the sub-measure with output flag 0 is $(1-e)\lambda\,\hat q_{0,0}(P)$, and the output class requires it to equal
$m'\lambda'$ (with $\lambda'=\lambda(p')$). Since $|\hat q_{0,0}(P)|\le s_0$,
\[
m'\lambda'\le(1-e)\lambda\,s_0.\tag{$*$}
\]
If $e'\ge e$ (that is, $m'\le1-e$): $m'\ge(1-e)s_0$, and substituting into $(*)$ gives $\lambda'\le\lambda$.
If $e'<e$: $s_0\le1$, and $(*)$ gives $(1-e')\lambda'\le(1-e)\lambda$.

The last sentence follows from Theorem~\ref{thm:monotone}.
\end{proof}

\paragraph{Level-1 form.}
For the uncoded qubit, (i) and (ii) are also equivalent to the existence of an arbitrary simulable reduction (label permutations beyond translations) and to $\Phi(P')\le\Phi(P)$ for all $\varphi\in\Phi_4$ (Proposition~\ref{prop:local-level1}); the complete family of monotones is described in Theorem~\ref{thm:complete}.

\begin{remark}[Structure]
(1) Complete monotones: the local free order is completely determined by two functionals: the total unflagged eigenvalue mass $\mu=(1-e)\lambda$ (nonincreasing under any local free operation)
and ``the conditional eigenvalue $\lambda$ is nonincreasing when the erasure rate does not decrease''.
(2) Slope of the boundary of the reachable set: the steepest direction for lowering $e'$ starting from $(p,e)$ is $\mu'=\mu$, whose tangent slope is exactly
$dp'/de'=-c(p',e')$, consistent with the $c$ of Theorem~\ref{thm:rate}: $c$ is the boundary slope of the free order.
(3) The order is not total; for example, if $(p,e)$ and $(p',e')$ satisfy $e'<e$ but $\mu'>\mu$, neither direction is reachable.
\end{remark}

\subsection{Hellinger/Chernoff monotones}

\begin{proposition}[Hellinger monotones]\label{prop:hellinger}
For a noise instance $P=P_{EF}$, $G\in V$, and $s\in(0,1)$ let
$\BC_s(P;G)=\sum_{E,F}P(E,F)^{1-s}P(E+G,F)^s$.
If $P'$ is obtained by a free operation of type (F1) or (F2)/(F2$'$), then $\BC_s(P';G)\ge\BC_s(P;G)$ (for the same $G$).\hfill\st{PROVED}
\end{proposition}
\begin{proof}
These operations are Markov kernels $W$ acting on $(E,F)$ (discarding/coarse-graining flags; superposition of an $E$-independent $N$ is a convolution),
and translation commutes with convolution: $(WP)(\cdot+G)=W(P(\cdot+G))$. The function $(x,y)\mapsto x^{1-s}y^s$ is positively homogeneous and concave, hence superadditive,
$\sum_u(\sum_vW(u|v)x_v)^{1-s}(\sum_vW(u|v)y_v)^s\ge\sum_v x_v^{1-s}y_v^s$. Take $x=P$, $y=P(\cdot+G)$.
\end{proof}

For $\Nn_{\rm cc}(p,e)$ and any nonzero single-qubit shift $G_j\in\{X,Y,Z\}$ (the three are symmetric), the single-qubit coefficient is
\[
b_s(p,e)=e+(1-e)\Big[(1-p)^{1-s}(\tfrac p3)^s+(\tfrac p3)^{1-s}(1-p)^s+\tfrac{2p}3\Big],
\]
where the first term comes from $F=1$ (the uniform distribution is invariant under the shift). Since $b_s=b_{1-s}$ and $b_s$ is convex in $s$, $\min_s b_s=b_{1/2}$. Put
\[
\beta_{\rm dep}(p)=2\sqrt{p(1-p)/3}+2p/3,\qquad B(p,e)=b_{1/2}(p,e)=e+(1-e)\beta_{\rm dep}(p).
\]
$B$ is the \emph{single-qubit Bhattacharyya parameter}. For each $s$, the same argument as in Theorem~\ref{thm:partial} together with Proposition~\ref{prop:hellinger}
(partial flag discarding does not decrease $b_s$) gives $\partial_eb_s\le c\,\partial_pb_s$, that is, the Hellinger exchange rate satisfies
\[
R_s(p,e)=\frac{\partial_eb_s}{\partial_pb_s}\le c(p,e).
\]
Write $R_B:=R_{1/2}=\dfrac{1-\beta_{\rm dep}(p)}{(1-e)\,\beta_{\rm dep}'(p)}$.

\subsection{The gap problem}\label{sec:gap}

The free order $\succeq_{\rm free}$ is a sub-order of the \emph{operational order} (comparing $\err_d$ or $\alpha$ for a given code). Let $R_\alpha$ be the true exchange rate on the surface code.
\begin{problem}[Exchange-rate gap]\label{prob:gap}
Quantify $\Delta:=c-R_\alpha\ge0$, and decide whether the gap comes from
\begin{enumerate}[(A)]
\item \emph{code dependence}: the \emph{code-universal} operational order (requiring $\err$ not to increase for all linear decoding structures and all $n$) coincides with the local free order,
and the $\alpha$-order of the surface code is merely a coarser order; or
\item \emph{incompleteness of the free operations}: there are free operations outside Definition~\ref{def:local} (nonlocal or code dependent) that would make the order finer?
\end{enumerate}
\end{problem}

\begin{conjecture}[Code-universal order = local free order]\label{conj:universal}
Within $\Nn_{\rm cc}(p,e)$, the \emph{code-universal finite-$n$ operational order} coincides with the local free order of Theorem~\ref{thm:preorder}.
\hfill\st{CONJECTURE}
\end{conjecture}
\noindent\emph{Reformulation (constitutional audit).} The all-losses version of this conjecture (compare $\Phi(P')\le\Phi(P)$ for all $\varphi\in\Phi_\Lambda$, over all structures) is proved by taking the uncoded qubit (Proposition~\ref{prop:local-level1}, Remark~\ref{rem:conj-reform}); what remains open is the single-loss statement for $\err=1-\Phi_{\max}$.

\noindent\emph{Evidence and route.} Observation~\ref{obs:sharp} gives first-order evidence at $e_0=0$: the boundary slope $c$ is attained by some code.
Route to a proof: (S1) for each $n$ and linear structure $(\sigma,L)$, the comparison of $\err$ is a Blackwell comparison of the ``experiments'' given by $n$ independent copies
(Blackwell--Sherman--Stein: experiment $A$ is at least as good as $B$ for all priors and losses if and only if $B$ is a randomized processing of $A$ \cite{Blackwell});
(S2) the decision problems arising from linear structures form a proper subclass of all decision problems, and one must show that they suffice to separate the two functionals $\mu$ and $(e,\lambda)$ of the local free order;
(S3) the asymptotic version is characterized, via the large-sample Blackwell order, by the family of all R\'enyi divergences \cite{MPST}, which corresponds exactly to $\{b_s\}_{s\in(0,1)}$ and its higher-order generalizations.
If the conjecture holds, the gap is case (A) and $c-R_\alpha$ measures information specific to the surface code rather than a deficiency of the free operations;
if it is refuted by a counterexample from some code (for example one other than $[[5,1,3]]$), new free operations are needed (case (B)).

\begin{conjecture}[Completeness of $B$: hypothesis $H_B$]\label{conj:HB}
In some neighborhood within $\Nn_{\rm cc}$, the $\alpha(p,e)$ of the surface code is a function of $B(p,e)$, so that $R_\alpha=R_B$.
\hfill\st{CONJECTURE}
\end{conjecture}
\noindent This is testable: the numerics of P1 give $R_\alpha$, to be compared with the analytic $R_B$. $R_\alpha\neq R_B$ would show that
the Bhattacharyya monotone is incomplete, and the next candidates are $b_s$ ($s\ne\tfrac12$). Lemma~\ref{lem:union} (upper bound) and Proposition~\ref{prop:genie}
(lower bound) already sandwich $\alpha$ between two functions depending only on $B$, but a sandwich does not constrain derivatives, so $H_B$ needs its own test.

\subsection{Analytic formula for derivatives and the envelope argument}\label{sec:envelope}

The sampling scheme of P1 (README \S2) obtains $\partial_\theta\pL$ directly from one table $f(k,w)$, without decoding at neighboring parameter points. For decoders independent of $(p,e)$ (MWPM) this is an identity;
for the maximum-likelihood decoder it relies on the following argument.

\begin{remark}[Envelope argument for ML derivatives]\label{rem:envelope}
Let $D^\ast$ be the ML decoder at $\theta^\ast=(p^\ast,e^\ast)$ (a fixed map). For a fixed decoder $D$,
$\pL(\theta;D)=\sum_{k,w}P_\theta(k,w)\,f_D(k,w)$, where $f_D$ does not depend on $\theta$. Optimality gives $\err(\theta)=\min_D\pL(\theta;D)\le\pL(\theta;D^\ast)$,
with equality at $\theta^\ast$. The difference $g(\theta)=\pL(\theta;D^\ast)-\err(\theta)\ge0$ satisfies $g(\theta^\ast)=0$, so wherever both are differentiable at $\theta^\ast$ we have $\nabla g(\theta^\ast)=0$, that is,
\[
\partial_\theta\err\big|_{\theta^\ast}=\sum_{k,w}\partial_\theta P_\theta(k,w)\,f_{D^\ast}(k,w).
\]
The ML decoder depends on $\theta$ only through $p$: the erased positions are known, so $e$ never enters the posterior given the flags. Therefore
the $e$-derivative is an exact identity for ML, and the envelope argument is needed only for $\partial_p$.
Hence $\partial_\theta\ln\pL=\langle s_\theta\rangle_{\rm fail}$ (with $s_\theta=\partial_\theta\ln P_\theta$ the score function and the conditional expectation taken over failure events),
and $\partial\alpha/\partial\theta$ is minus the least-squares slope of $\langle s_\theta\rangle_{\rm fail}$ against $d$ (the least-squares slope is linear in the data and commutes with differentiation).
Differentiability: at finite $d$, $\err$ is the minimum of finitely many smooth functions and is differentiable almost everywhere; kinks may occur at ties between decisions.
\hfill\st{SKETCH}
\end{remark}

\begin{observation}[Exact verification of the envelope identity]\label{obs:envelope}
Check C8 of \texttt{exact\_small\_codes.py}: for $[[5,1,3]]$, the Steane code, and the $d=3$ rotated surface code, at $(p^\ast,e^\ast)\in\{(0.04,0.05),(0.1,0.2)\}$,
the derivatives $\partial_p\err,\partial_e\err$ of the exact ML agree with the derivatives of the failure rate of the \emph{frozen decoder} to within relative $2\times10^{-10}$.
In addition, \texttt{experiments/p1\_erasure\_pauli/envelope\_check.py} tests the $p$-derivative at $d=5,7$ with the exact transfer-matrix ML decoder on common random samples
of each stratum $(k,w)$. A central difference has a deterministic $O(h^2)$ discretisation error (about $+0.3\%$ at $d=5$ and $+1.9\%$ at $d=7$ for $h=0.01$, $p=0.06$, growing four times for $2h$),
so the test compares the finite difference with re-matched decoders ($B_{\rm matched}$) against the \emph{same} finite difference with the frozen decoder ($B_{\rm frozen}$).
Their difference is at most $0.04\%$ of the derivative (within $1.6$ bootstrap standard errors) in four cases, one of which has erasure strata; one deviation of $2.05$ standard errors ($0.1\%$ of the derivative) occurs at step $2h$, $d=7$.
The identity has not yet been tested directly for the truncated-MPS ML decoder or at $d\ge9$; the MPS decoder's decisions agree with the exact decoder on all tested samples (README, M1), so it is expected to transfer.\hfill\st{NUMERICAL}
\end{observation}
